%%=============================================================================
%% Voorwoord
%%=============================================================================

\chapter*{\IfLanguageName{dutch}{Woord vooraf}{Preface}}%
\label{ch:voorwoord}

%% TODO:
%% Het voorwoord is het enige deel van de graduaatsproef waar je vanuit je
%% eigen standpunt (``ik-vorm'') mag schrijven. Je kan hier bv. motiveren
%% waarom jij het onderwerp wil bespreken.
%% Vergeet ook niet te bedanken wie je geholpen/gesteund/... heeft

Deze graduaatsproef is voor mij een kans geweest om de leerstof van de opleiding Graduaat Programmeren toe te passen in een project dat dicht bij een realistische werksituatie ligt.
Met Forcebit Ticketing wilde ik de geziene leerstof herhalen en combineren in een oplossing voor een concreet probleem bij Forcebit: communicatie met klanten stond verspreid over telefoon en e-mail, waardoor opvolging minder overzichtelijk was.
De applicatie brengt tickets, gesprekken, statussen en bijlagen samen in één centrale workflow.

\medskip

Technisch bestaat het project uit een gelaagde backend in C\# en ASP.NET Core, gekoppeld aan een MySQL-databank.
Daarin gebruikte ik onder andere Dependency Injection, het Repository Pattern, een Service Layer, DTO-mapping en het Options Pattern.
Ook de frontend is gestructureerd opgebouwd met aparte schermen, custom hooks, API-modules, formulieren, componenten, validatieschema's en utility-functies, zodat de code herbruikbaar en beter uitlegbaar blijft.

\medskip

Tijdens dit project heb ik ook nieuwe onderdelen leren toepassen, zoals een mailservice, het uploaden van attachments, previews voor afbeeldingen, authenticatie met rollen en een duidelijke ticketworkflow.
Daardoor voelde deze graduaatsproef voor mij als een passende afsluiter van de opleiding: een herhaling van wat ik geleerd heb, aangevuld met technieken die dichter bij echte softwareontwikkeling liggen.

\medskip

Graag wil ik mijn mentor Emlin Oguz Inci bedanken voor het opvolgen van deze graduaatsproef en voor zijn feedback over de features, structuur en verdedigbaarheid van het project.
Zijn opmerkingen hebben mij geholpen om niet alleen naar werkende code te kijken, maar ook naar gebruiksgemak, onderhoudbaarheid en duidelijke uitleg.

